% sections/abstract.tex -- the abstract (paper_v2_vn, translated in L2 from paper_v2's
% sections/abstract.tex, as the sample of the style; terms as in glossary_vn.tex). Every number is
% \val{} from numbers.tex.
HAN và các mô hình tương tự khai thác đồ thị không đồng nhất thông qua các meta-path, vốn thường
được chọn thủ công. Các meta-path do phương pháp tìm kiếm LMSPS tìm được đã giúp HAN phân loại nút
tốt hơn, nhưng việc tìm kiếm như vậy có đáng với chi phí trong bài toán gợi ý hay không thì chưa
được kiểm chứng. Trên Amazon-Meta, chúng tôi so sánh hai cách chọn meta-path cho HAN từ cùng một tập
\val{search-cands} ứng viên: thủ công (cố định từ trước) và bằng LMSPS được điều chỉnh cho bài toán
dự đoán liên kết (tổng hợp từ \val{search-runs} lần chạy), với các tập chọn ngẫu nhiên làm đối
chứng. Mọi tập được thử nghiệm trên HAN và metapath2vec. Các tập tìm được kém các tập thủ công
\val{han-6-hit-relabs}--\val{han-4-hit-relabs}\% theo độ đo đã cố định từ trước và chỉ ngang bằng
khi xếp hạng toàn bộ, sau khoảng \val{search-total-h}~giờ~GPU tìm kiếm. Phần thua kém chủ yếu nằm
ở sách và nhạc, những sản phẩm mà các meta-path danh mục của các tập tìm được không đến được. Đối
chứng ngẫu nhiên cho thấy việc chọn meta-path ít ảnh hưởng đến HAN trên lược đồ này: mọi tập cùng
kích thước chỉ chênh nhau trong phạm vi \val{han-4-span} Hit@10 và các tập chọn ngẫu nhiên đạt kết
quả ít nhất ngang bằng cả hai cách. Với HAN trên lược đồ nhỏ này, kết quả của tìm kiếm tự động khá
gần với thiết kế thủ công, vì vậy lựa chọn giữa hai cách tùy thuộc vào bối cảnh: thiết kế thủ công
tốt hơn khi có chuyên gia làm được việc này trong thời gian ngắn hơn thời gian tìm kiếm, còn tìm
kiếm cho các tập gần tương đương mà không cần kiến thức chuyên môn.
